\section{Introduction}

Vision-based tactile sensors such as GelSight, GelSlim and DIGIT~\cite{yuan2017gelsight,donlon2018gelslim,lambeta2020digit} measure contact by imaging a soft elastomer under internal colored illumination. The image is not a picture of the object: it is the result of gel deformation, LED placement and color, elastomer reflectance, camera response and manufacturing variation. Two sensors of the same design, or the same sensor across sessions, produce different images for the same contact. Dense perception from these images---contact depth, surface normals, object pose---therefore needs sensor-specific supervision, and dense real labels are expensive to obtain.

Real-to-sim-to-real pipelines address this by \emph{calibrating a tactile renderer} on a small amount of real data and then training on synthetic images whose depth, normal and pose labels come for free~\cite{si2022taxim,wang2022tacto,chen2023tacchi,zhong2025tactgen}. Two renderer families are in use. \emph{Physically calibrated} renderers keep a light-transport model of the sensor and fit its parameters---LED color and strength, gel material, camera response---to real images~\cite{agarwal2021simulation,si2022taxim}. \emph{Learned} renderers replace the light-transport model by an image-to-image network trained on real (depth, image) pairs~\cite{zhong2025tactgen}. The two families sit at opposite ends of a data-efficiency axis: a physical renderer can be calibrated from a single no-contact frame, whereas a learned renderer needs labeled contact pairs, i.e. exactly the data the pipeline is meant to save.

This paper asks how much real data each renderer needs before its synthetic data improves a downstream tactile model, and whether the reported gains survive a control that is usually missing: the renderer must not see the real frames on which the downstream model is tested. In many learned-renderer pipelines the renderer is trained on all available real pairs and the downstream model is validated on frames from the same sessions; the renderer has then memorized the appearance of the test contacts, and the synthetic data carries that information into training. We show this \emph{leakage} inflates the measured benefit by up to a factor of two.

We study the question on \sensor{}, a GelSlim-type fingertip sensor, with a single downstream architecture (a frozen sensor-invariant tactile encoder~\cite{gupta2025sitr} with dense depth/normal and planar-pose heads) and two renderers built on the same object meshes and contact sampler: a Blender renderer calibrated by Bayesian optimization to one no-contact image, and a pix2pix depth-to-image renderer~\cite{isola2017pix2pix,church2022tactile} trained on $\Kreal$ real pairs per object. We evaluate under two protocols. \emph{Leave-object-out} holds out three objects with no real data at all and asks whether simulated contacts of their meshes transfer. \emph{$\Kreal$-shot} gives every object $\Kreal\in\{0,10,25,100\}$ real frames, and---critically---trains the learned renderer on the same $\Kreal$ frames as the downstream model.

Our findings are:
\begin{enumerate}
    \item \textbf{Simulation is decisive for pose transfer to new objects.} With no real data of an object, a tactile pose model fails ($48^\circ$ mean rotation error). Adding simulated contacts of the object's mesh reduces this to $5.6^\circ$ with the physically calibrated renderer and $6.3^\circ$ with the strict learned renderer. A control that keeps the renderer but removes the held-out geometry from simulation stays at $55^\circ$: the benefit is geometric coverage, and the renderer's appearance fidelity is secondary.
    \item \textbf{For dense depth, the benefit is gated by appearance fidelity and by $\Kreal$.} The physically calibrated renderer does not lower real depth error at any $\Kreal$. The learned renderer lowers it only from $\Kreal=100$ pairs per object ($-18\%$ MSE, $-15\%$ contact-region MAE, $+0.04$ contact IoU); at $\Kreal\le 25$ it is indistinguishable from training on the real frames alone. Few-shot depth error is dominated by overfitting to the few real frames, not by the amount of synthetic data.
    \item \textbf{Leakage matters.} A learned renderer trained on all real pairs appears to cut depth error by $24$--$31\%$ at every $\Kreal$ and on unseen objects; under the strict protocol the gain is $0$--$18\%$. We recommend reporting learned-renderer results with the renderer's training frames disjoint from the downstream test frames.
    \item \textbf{A cheap fidelity proxy predicts downstream ranking.} The contact-region L1 between rendered and real images on held-out pairs orders eight renderer variants the same way as their downstream depth error, so renderer design can be iterated in minutes rather than hours.
    \item \textbf{Negative results on the appearance side.} Pretraining the learned renderer on the physical renderer's images, normalizing per-session illumination drift (which we measure to be larger than the renderer error itself), and predicting difference images all fail to improve small-$\Kreal$ fidelity in the contact region, locating the bottleneck in contact shading rather than in the illumination field.
\end{enumerate}
